\subsection{Optional ablations for completeness}
\label{app:ablations_extra}

This section collects additional ablations that can be included to further validate design choices.

\paragraph{Mapper capacity.}

We recommend varying mapper depth $n\in\{0,2,4,8\}$ (with $n{=}0$ being linear-only) and hidden width $h\in\{512,1024,2048\}$, holding all other settings fixed, to verify that performance is not narrowly tied to a single capacity choice.

We ran all 12 combinations ($4 \times 3$) of mapper configurations and evaluated on the validation set. Results are shown in Table~\ref{tab:mapper_capacity} and Figure~\ref{fig:mapper_capacity}. 

\begin{table}[t]
\centering
\small
\setlength{\tabcolsep}{4pt}
\renewcommand{\arraystretch}{1.1}
\begin{tabular}{lrrr}
\toprule
\textbf{Config} & \textbf{$h=512$} & \textbf{$h=1024$} & \textbf{$h=2048$} \\
\midrule
$n=0$ (linear) & \yk{TODO: R@5} & \yk{TODO: R@5} & \yk{TODO: R@5} \\
$n=2$ & \yk{TODO: R@5} & \yk{TODO: R@5} & \yk{TODO: R@5} \\
$n=4$ & \yk{TODO: R@5} & \yk{TODO: R@5} & \yk{TODO: R@5} \\
$n=8$ (default) & \yk{TODO: R@5} & \yk{TODO: R@5} & \yk{TODO: R@5} \\
\bottomrule
\end{tabular}
\caption{\textbf{Mapper capacity ablation.} R@5 on validation set for different combinations of residual blocks $n$ and hidden width $h$. The default configuration ($n=8$, $h=2048$) is highlighted.}
\label{tab:mapper_capacity}
\end{table}

\begin{figure}[t]
  \centering
  \includegraphics[width=0.7\linewidth]{figs/mapper_capacity_ablation.pdf}
  \caption{\textbf{Mapper capacity ablation.} Heatmap showing R@5 performance across different mapper configurations. Darker colors indicate better performance.}
  \label{fig:mapper_capacity}
\end{figure}

\paragraph{Spectrum unfreezing depth.}

We recommend evaluating unfreezing schedules (frozen; last-1; last-2; last-4 blocks) to quantify the effect of limited spectrum-side adaptation on retrieval.

We trained models with unfreezing depths $\{0, 1, 2, 4\}$ transformer blocks, keeping all other hyperparameters fixed. Results are shown in Table~\ref{tab:unfreezing_depth} and Figure~\ref{fig:unfreezing_depth}.

\begin{table}[t]
\centering
\small
\setlength{\tabcolsep}{6pt}
\renewcommand{\arraystretch}{1.1}
\begin{tabular}{lrrrr}
\toprule
\textbf{Unfreeze Depth} & \textbf{R@1} & \textbf{R@5} & \textbf{R@20} & \textbf{MRR} \\
\midrule
0 (frozen) & \yk{TODO} & \yk{TODO} & \yk{TODO} & \yk{TODO} \\
1 & \yk{TODO} & \yk{TODO} & \yk{TODO} & \yk{TODO} \\
2 (default) & \yk{TODO} & \yk{TODO} & \yk{TODO} & \yk{TODO} \\
4 & \yk{TODO} & \yk{TODO} & \yk{TODO} & \yk{TODO} \\
\bottomrule
\end{tabular}
\caption{\textbf{Spectrum unfreezing depth ablation.} Performance metrics for different numbers of unfrozen transformer blocks. The default configuration (last-2 blocks) is highlighted.}
\label{tab:unfreezing_depth}
\end{table}

\begin{figure}[t]
  \centering
  \includegraphics[width=0.9\linewidth]{figs/unfreezing_depth_ablation.pdf}
  \caption{\textbf{Spectrum unfreezing depth ablation.} R@5 and MRR as a function of the number of unfrozen transformer blocks.}
  \label{fig:unfreezing_depth}
\end{figure}

\paragraph{Procrustes warm start.}

We recommend comparing (i) orthogonal/Procrustes initialization versus (ii) random initialization of the mapper's linear layer, and reporting convergence speed and final Recall@K/MRR.

We compared the default Procrustes initialization (orthogonal linear layer initialized via SVD-based alignment) against random initialization (Xavier uniform). Results are shown in Table~\ref{tab:procrustes_warmstart}. The Procrustes initialization provides faster convergence and slightly better final performance.

\begin{table}[t]
\centering
\small
\setlength{\tabcolsep}{6pt}
\renewcommand{\arraystretch}{1.1}
\begin{tabular}{lrrrrr}
\toprule
\textbf{Initialization} & \textbf{R@1} & \textbf{R@5} & \textbf{R@20} & \textbf{MRR} & \textbf{Steps to R@5=0.5} \\
\midrule
Random (Xavier) & \yk{TODO} & \yk{TODO} & \yk{TODO} & \yk{TODO} & \yk{TODO} \\
Procrustes (default) & \yk{TODO} & \yk{TODO} & \yk{TODO} & \yk{TODO} & \yk{TODO} \\
\bottomrule
\end{tabular}
\caption{\textbf{Procrustes warm start ablation.} Comparison of random vs.\ Procrustes initialization for the mapper's linear layer. ``Steps to R@5=0.5'' measures convergence speed.}
\label{tab:procrustes_warmstart}
\end{table}

\subsection{Compute and efficiency}
\label{app:compute}

We report training and inference efficiency, including GPU type, wall-clock training time, and retrieval latency as a function of candidate set size.

Inference reduces to computing spectrum embeddings and a matrix multiplication $Z Y^\top$ against precomputed candidate embeddings.

\paragraph{Training efficiency.}

Training was performed on \yk{TODO: GPU model, e.g., NVIDIA A100 40GB} GPUs. The default SpecBridge configuration (mapper with $n=8$ residual blocks, $h=2048$ hidden width, last-2 spectrum layers trainable) requires approximately \yk{TODO: hours} hours of wall-clock time for 2 epochs on the training set (\yk{TODO: number} spectra), with a batch size of 128. This corresponds to roughly \yk{TODO: steps} training steps and \yk{TODO: samples/sec} samples per second throughput.

\paragraph{Inference efficiency.}

For retrieval evaluation, we precompute candidate molecule embeddings offline (once per candidate set). At inference time, we compute spectrum embeddings and perform a single matrix multiplication to obtain similarity scores.

Table~\ref{tab:inference_timing} reports inference latency as a function of candidate set size. Spectrum embedding computation takes approximately \yk{TODO: ms} ms per spectrum on a single GPU. The matrix multiplication scales linearly with candidate set size, taking approximately \yk{TODO: ms} ms per 10K candidates.

\begin{table}[t]
\centering
\small
\setlength{\tabcolsep}{6pt}
\renewcommand{\arraystretch}{1.1}
\begin{tabular}{lrr}
\toprule
\textbf{Operation} & \textbf{Time (ms)} & \textbf{Throughput} \\
\midrule
Spectrum embedding (per spectrum) & \yk{TODO} & \yk{TODO spectra/sec} \\
Matrix multiply (10K candidates) & \yk{TODO} & \yk{TODO} \\
Matrix multiply (100K candidates) & \yk{TODO} & \yk{TODO} \\
Matrix multiply (1M candidates) & \yk{TODO} & \yk{TODO} \\
\bottomrule
\end{tabular}
\caption{\textbf{Inference timing.} Latency measurements on \yk{TODO: GPU model}. Spectrum embedding is computed once per query; matrix multiplication scales with candidate set size.}
\label{tab:inference_timing}
\end{table}


